\documentclass[11pt,a4paper]{article}
\usepackage[margin=1in]{geometry}
\usepackage{hyperref}
\usepackage{amsmath}
\usepackage{enumitem}
\usepackage{graphicx}
\usepackage{xcolor}
\usepackage{tabularx}
\usepackage{booktabs}
\usepackage{array}
\usepackage{longtable}
\usepackage{titling}

% -----------------------------------------------------
% UCS503 Project Proposal
% -----------------------------------------------------

\definecolor{BattifyBlue}{HTML}{1E3A5F}
\definecolor{BattifyAccent}{HTML}{D97706}

\newcommand{\BvrHint}[1]{{\normalfont\normalsize\itshape\hfill #1}}

\title{\textbf{BATTIFY}\\
\large On-Demand Battery Delivery, Installation and Buyback Platform}

\pretitle{\begin{center}\bfseries\fontsize{18bp}{22bp}\selectfont}
\posttitle{\par\end{center}\vskip 0.5em}

\author{
Vyomesh Joshi, CSED\\
Ph No: \texttt{6284329484}\\
Roll No: \texttt{1024030181}\\
Email: \texttt{vjoshi\_be24@thapar.edu}\\[4pt]
Sarika Garg, CSED\\
Ph No: \texttt{6284532323}\\
Roll No: \texttt{1024030741}\\
Email: \texttt{sgarg1\_be24@thapar.edu}\\[4pt]
Shivinder Singh, CSED\\
Ph No: \texttt{9653797001}\\
Roll No: \texttt{1024030213}\\
Email: \texttt{ssingh25\_be24@thapar.edu}\\[4pt]
Yatish Goyal, CSED\\
Ph No: \texttt{9888327011}\\
Roll No: \texttt{1024030187}\\
Email: \texttt{ygoyal\_be24@thapar.edu}
}

\date{\today}

\begin{document}
\maketitle

\begin{center}
\textbf{\color{BattifyAccent} Project Proposal for UCS503P -- Software Engineering Lab}\\
Thapar Institute of Engineering and Technology
\end{center}

\section{Higher-order Goal \BvrHint{Real-world impact}}

Battify aims to make battery replacement and servicing a fast, reliable and
transparent service rather than a conventional retail purchase. The platform
targets households, farmers, vehicle owners and small businesses that may need
a replacement battery urgently but currently have to locate a retailer,
arrange transportation, wait for availability and separately handle
installation and disposal of the old battery.

The engineering objective is to build a software system that coordinates
customers, inventory, delivery personnel, installation and old-battery
buyback through one workflow. The project therefore has a direct real-world
use case while providing enough software complexity to demonstrate the
learning outcomes of UCS503.

\section{Problem Statement}

Battery failure is often time-sensitive. A failed inverter battery can leave a
household without backup power, while a failed tractor or agricultural pump
battery can interrupt farm operations. Existing purchasing methods commonly
separate the activities of finding a battery, checking availability,
transporting it, installing it and disposing of the old unit.

The major problems are:

\begin{itemize}[leftmargin=0.7cm]
    \item \textbf{Fragmented service:} battery purchase, delivery, installation
    and old-battery disposal are handled separately.
    \item \textbf{Uncertain availability:} customers may not know whether the
    required battery is actually in stock nearby.
    \item \textbf{Slow replacement:} urgent customers have limited options for
    same-day delivery and installation.
    \item \textbf{Poor workflow visibility:} customers often cannot track
    whether an order has been accepted, dispatched or installed.
    \item \textbf{Seasonal demand:} agricultural batteries experience demand
    spikes around sowing and harvesting periods, making inventory planning
    important.
    \item \textbf{Battery waste:} customers need a convenient mechanism for
    returning old batteries through a responsible recovery channel.
\end{itemize}

\section{Proposed Solution \BvrHint{Software Proposition}}

Battify will be developed as a web-based platform connecting customers,
inventory managers and delivery/installation agents.

The proposed system will provide:

\begin{itemize}[leftmargin=0.7cm]
    \item \textbf{Battery catalogue:} customers can select battery type,
    capacity, vehicle/application and preferred brand.
    \item \textbf{Location-aware ordering:} the system records the delivery
    location and identifies suitable nearby inventory.
    \item \textbf{Inventory management:} administrators can monitor stock,
    reserved units, low-stock products and seasonal demand.
    \item \textbf{Delivery and installation workflow:} orders move through
    defined states from placement to installation.
    \item \textbf{Old-battery buyback:} the old battery is recorded and its
    estimated buyback value is deducted from the order.
    \item \textbf{Customer tracking:} customers can see order status and
    installation progress.
    \item \textbf{Operations dashboard:} administrators can monitor orders,
    delivery times, stock levels and unresolved jobs.
\end{itemize}

\section{Core Workflow}

The principal workflow is:

\begin{enumerate}[leftmargin=0.7cm]
    \item Customer selects the required battery and provides vehicle or
    application details.
    \item The system validates the request and checks inventory.
    \item The customer selects a delivery location and preferred service slot.
    \item An order is created and assigned an order ID.
    \item The operations dashboard assigns the order to a delivery/installation
    agent.
    \item The agent collects or verifies the old battery, delivers the new
    battery and performs installation.
    \item The old battery is recorded for buyback/recycling.
    \item The order is marked completed and the customer receives a final
    status and invoice.
\end{enumerate}

\subsection{Order State Model}

A controlled state machine can be used:

\[
\text{Placed}
\rightarrow
\text{Confirmed}
\rightarrow
\text{Assigned}
\rightarrow
\text{Dispatched}
\rightarrow
\text{Delivered}
\rightarrow
\text{Installed}
\rightarrow
\text{Completed}
\]

Exceptional states such as \textit{Cancelled}, \textit{Out of Stock} and
\textit{Rescheduled} will also be supported.

\section{Why This is Suitable for UCS503}

Battify is intentionally designed as more than a simple e-commerce website.
The core challenge is the coordination of multiple actors, inventory,
geographical service, order states, installation and reverse logistics.

This creates a practical software engineering problem involving:

\begin{itemize}[leftmargin=0.7cm]
    \item requirements elicitation and validation;
    \item UML and workflow modelling;
    \item modular system and database design;
    \item architecture and API design;
    \item software testing and verification;
    \item Agile backlog and sprint management;
    \item energy-aware software and UI decisions;
    \item project scope, time and cost estimation.
\end{itemize}

\section{Mapping to Course Learning Outcomes}

\begin{longtable}{p{0.19\textwidth}p{0.73\textwidth}}
\toprule
\textbf{CLO} & \textbf{How Battify Demonstrates the CLO} \\
\midrule
\textbf{CLO 1} &
A software process model can be selected and justified for the project.
An incremental/Agile approach is suitable because the system can be delivered
in modules such as catalogue, ordering, inventory, dispatch and buyback.
The project can compare Waterfall, Incremental and Agile approaches before
selecting the final process. \\

\textbf{CLO 2} &
Requirements can be elicited from multiple stakeholders: customers,
administrators, inventory managers, delivery agents and installation
personnel. Use-case, activity, sequence, class and state diagrams can model
the resulting requirements, followed by a formal SRS. \\

\textbf{CLO 3} &
The platform naturally supports separation of concerns between customer UI,
operations dashboard, order management, inventory, dispatch and buyback.
Database design, API design, MVC architecture, modularity, coupling and
cohesion can all be demonstrated. \\

\textbf{CLO 4} &
The order lifecycle provides clear testable states and business rules.
Unit tests can cover price calculation, inventory reservation and buyback
calculation; integration tests can cover order-to-dispatch workflows; system
tests can cover complete customer journeys. Black-box and white-box testing
strategies can also be documented. \\

\textbf{CLO 5} &
The project can be managed through Agile user stories, product backlog,
sprint backlog, story points, velocity and sprint reviews. UI sketches and
story cards can be used before implementation, while Test Driven Development
can be applied to critical modules. \\

\bottomrule
\end{longtable}

\section{Requirements Engineering}

\subsection{Stakeholders}

The principal stakeholders are:

\begin{itemize}[leftmargin=0.7cm]
    \item Customer
    \item Delivery and installation agent
    \item Inventory manager
    \item Operations administrator
    \item Battery supplier
    \item Buyback/recycling partner
\end{itemize}

\subsection{Representative Functional Requirements}

\begin{enumerate}[leftmargin=0.7cm]
    \item The system shall allow customers to browse battery products.
    \item The system shall allow customers to submit a replacement request.
    \item The system shall check inventory before confirming an order.
    \item The system shall reserve inventory for confirmed orders.
    \item The system shall assign an eligible service agent to an order.
    \item The system shall maintain an auditable order-status history.
    \item The system shall calculate the final amount after applicable buyback
    credit.
    \item The system shall allow administrators to update inventory.
    \item The system shall provide customers with order tracking.
    \item The system shall record completed installations and old-battery
    collection.
\end{enumerate}

\subsection{Representative Non-functional Requirements}

\begin{itemize}[leftmargin=0.7cm]
    \item \textbf{Performance:} common customer operations should respond
    quickly under expected project-scale load.
    \item \textbf{Usability:} the customer workflow should require minimal
    steps and work well on mobile devices.
    \item \textbf{Reliability:} order and inventory state changes should be
    consistent and recoverable.
    \item \textbf{Security:} authentication and role-based authorization
    should protect administrative functions.
    \item \textbf{Maintainability:} the backend should use modular components
    with clear interfaces.
    \item \textbf{Energy efficiency:} the UI should avoid unnecessary
    animations, excessive polling and resource-heavy media.
\end{itemize}

\section{UML and System Modelling}

The project will produce the following models:

\begin{itemize}[leftmargin=0.7cm]
    \item \textbf{Use-case diagram:} customer, agent and administrator
    interactions.
    \item \textbf{Activity diagram:} battery ordering and installation flow.
    \item \textbf{Sequence diagram:} order creation, inventory validation and
    agent assignment.
    \item \textbf{Class diagram:} Customer, Order, Battery, InventoryItem,
    Agent, Installation and Buyback entities.
    \item \textbf{State diagram:} order lifecycle and exceptional states.
    \item \textbf{ER diagram:} relational data model for users, products,
    inventory, orders, payments and service records.
    \item \textbf{Data Flow Diagram:} movement of information between
    customers, system modules, administrators and service agents.
\end{itemize}

\section{Software Architecture}

A three-tier architecture with MVC principles is proposed.

\begin{itemize}[leftmargin=0.7cm]
    \item \textbf{Presentation layer:} responsive web interface for customers
    and staff.
    \item \textbf{Application layer:} REST APIs and business logic for
    authentication, ordering, inventory, dispatch and buyback.
    \item \textbf{Data layer:} relational database containing users, products,
    inventory, orders, status history and service records.
\end{itemize}

The design will initially use a modular monolithic architecture rather than
microservices. This keeps the educational project manageable while retaining
clear module boundaries. The design can document how individual modules could
later be separated if operational scale requires it.

\section{Green Software Engineering}

Green software is explicitly relevant to Battify because the course includes
green coding, green architecture and energy-aware UI/UX.

The project will incorporate:

\begin{itemize}[leftmargin=0.7cm]
    \item lightweight pages and optimized images;
    \item reduced animations and unnecessary visual effects;
    \item pagination instead of loading entire datasets;
    \item indexed database queries for frequently accessed data;
    \item caching of relatively static catalogue information;
    \item event-driven updates where practical instead of aggressive polling;
    \item efficient API payloads;
    \item dark-mode support and reduced-power UI patterns;
    \item monitoring of unnecessary CPU, memory and network consumption.
\end{itemize}

For the course report, these decisions can be evaluated as engineering
trade-offs rather than simply being presented as visual features.

\section{Software Design Principles}

The system will demonstrate:

\begin{itemize}[leftmargin=0.7cm]
    \item \textbf{Separation of concerns:} UI, business logic and persistence
    remain independently organized.
    \item \textbf{Information hiding:} database implementation details remain
    behind service/repository interfaces.
    \item \textbf{High cohesion:} inventory operations are grouped within the
    inventory module, while dispatch operations remain in the dispatch module.
    \item \textbf{Low coupling:} modules communicate through well-defined APIs
    and interfaces.
    \item \textbf{Abstraction:} common operations such as order creation and
    status transitions are represented through service-level abstractions.
\end{itemize}

\section{Testing and Verification}

Testing will be planned alongside development.

\subsection{Unit Testing}

Examples include:

\begin{itemize}[leftmargin=0.7cm]
    \item buyback-credit calculation;
    \item total-price calculation;
    \item inventory availability checks;
    \item order-state transition validation;
    \item agent eligibility logic.
\end{itemize}

\subsection{Integration Testing}

Examples include:

\begin{itemize}[leftmargin=0.7cm]
    \item customer order $\rightarrow$ inventory reservation;
    \item inventory reservation $\rightarrow$ agent assignment;
    \item completed installation $\rightarrow$ inventory update;
    \item old-battery collection $\rightarrow$ buyback record.
\end{itemize}

\subsection{System and Acceptance Testing}

A complete customer scenario will be tested from battery selection through
order confirmation, delivery, installation, buyback and completion.

Black-box test cases will validate expected user-visible behavior, while
white-box testing will target important branches and business rules.

\section{Agile Development Plan}

The project can be organized into short sprints.

\begin{center}
\begin{tabularx}{\textwidth}{p{0.17\textwidth}X}
\toprule
\textbf{Sprint} & \textbf{Deliverable} \\
\midrule
Sprint 1 & Requirements, stakeholder analysis, SRS outline and UI sketches. \\
Sprint 2 & Authentication, customer catalogue and battery selection. \\
Sprint 3 & Order creation, inventory and pricing modules. \\
Sprint 4 & Agent assignment, delivery tracking and installation workflow. \\
Sprint 5 & Buyback module, admin dashboard and notifications. \\
Sprint 6 & Testing, debugging, green-software optimization and deployment. \\
\bottomrule
\end{tabularx}
\end{center}

A product backlog will contain user stories such as:

\begin{quote}
As a customer, I want to see available batteries for my vehicle so that I
can select a compatible replacement.
\end{quote}

\begin{quote}
As an operations manager, I want to see low-stock batteries so that I can
replenish inventory before demand causes service failures.
\end{quote}

\begin{quote}
As a delivery agent, I want to see my assigned installations so that I can
complete deliveries efficiently.
\end{quote}

\section{Evaluation Criteria}

The system will be evaluated using measurable software and operational
metrics:

\begin{itemize}[leftmargin=0.7cm]
    \item percentage of successful order submissions;
    \item average API response time for core operations;
    \item inventory consistency after concurrent order attempts;
    \item percentage of test cases passed;
    \item order completion rate in the simulated workflow;
    \item average number of steps required to place an order;
    \item reduction in unnecessary network requests after green-UI
    optimization;
    \item usability feedback from pilot users.
\end{itemize}

\section{Real-world Impact}

Battify is intended to solve a real operational problem rather than being
only a demonstration CRUD application. The platform can reduce the friction
between battery failure and replacement by integrating discovery, stock
verification, delivery, installation and old-battery recovery.

The buyback component also gives the system a reverse-logistics dimension.
Instead of treating the old battery as waste at the customer level, the
software records its collection and creates a traceable handoff to the
recovery/recycling process.

The agricultural use case is particularly relevant to Punjab because the
system can model seasonal inventory planning for tractor and pump batteries.
This gives the project an opportunity to demonstrate demand-aware software
rather than only static product ordering.

\section{Project Scope}

\subsection{Minimum Viable Product}

The MVP will contain:

\begin{itemize}[leftmargin=0.7cm]
    \item customer registration/login;
    \item battery catalogue;
    \item compatibility and requirement form;
    \item inventory database;
    \item order creation;
    \item buyback-credit calculation;
    \item administrator dashboard;
    \item delivery-agent assignment;
    \item order-status timeline;
    \item installation completion record;
    \item unit and integration test suite.
\end{itemize}

\subsection{Future Extensions}

Possible extensions include:

\begin{itemize}[leftmargin=0.7cm]
    \item WhatsApp-based order intake;
    \item route optimization for delivery agents;
    \item demand forecasting for seasonal agricultural batteries;
    \item automated compatibility recommendations;
    \item supplier-side inventory integration;
    \item predictive battery replacement reminders;
    \item multi-city expansion;
    \item dedicated mobile application.
\end{itemize}

\section{Risks and Mitigations}

\begin{itemize}[leftmargin=0.7cm]
    \item \textbf{Scope creep:} restrict the assessed implementation to the
    MVP and treat advanced features as future extensions.
    \item \textbf{Battery compatibility errors:} require structured vehicle,
    application and battery specifications rather than relying only on free
    text.
    \item \textbf{Inventory inconsistency:} use transactions and explicit
    reservation states.
    \item \textbf{Privacy and security:} collect only required customer
    information and implement role-based access.
    \item \textbf{Complex logistics:} simulate dispatch operations initially
    instead of attempting a full commercial logistics network.
\end{itemize}

\section{Summary}

Battify is proposed as an on-demand battery delivery, installation and
buyback platform that combines e-commerce, inventory management, field
service and reverse logistics.

Its strongest value for UCS503 is that the project naturally covers the
complete software engineering lifecycle: requirements engineering, UML
modelling, SRS documentation, architecture and design, coding practices,
testing, Agile development, project management and green software practices.

The existing Battify website can serve as the initial product prototype,
while the UCS503 lab implementation will focus on turning the concept into a
well-engineered, testable and documented software system.

\end{document}
